\par\Needspace{4\baselineskip}
\begin{center}
\SourceLeaf{023b}\pdfbookmark[2]{5.}{article-005}5.
\end{center}
Der Winkel $v$ liegt bei der Parabel zwischen $-180^0$ und $+180^0$, bei der Hyperbel zwischen $-(180^0-\psi)$ und $+(180^0-\psi)$. Bei der Ellipse dagegen durchläuft er in immer wiederkehrenden Perioden den ganzen Kreis. Wir nennen ihn die \emph{wahre Anomalie} des bewegten Körpers. Bisher haben allerdings fast alle Astronomen die wahre Anomalie bei der Ellipse vom Aphel statt vom Perihel an gezählt, entgegen der Analogie zur Parabel und zur Hyperbel, bei denen es kein Aphel gibt und man daher vom Perihel ausgehen musste. Bei der Wiederherstellung der Analogie zwischen allen Arten von Kegelschnitten
\SourceLeaf{024}zögerten wir umso weniger, als uns die französischen Astronomen der jüngsten Zeit bereits mit ihrem Beispiel vorangegangen waren.

Oft ist es zweckmäßig, die Form des Ausdrucks $r = \frac{p}{1+e\cos v}$ etwas zu ändern. Insbesondere sind die folgenden Formen zu beachten.
\[
r = \frac{p}{1+e-2e\sin\frac{1}{2}v^2} = \frac{p}{1-e+2e\cos\frac{1}{2}v^2}
\]
\[
r = \frac{p}{(1+e)\cos\frac{1}{2}v^2+(1-e)\sin\frac{1}{2}v^2}\cdot
\]

Für die Parabel gilt also $r = \frac{p}{2\cos\frac{1}{2}v^2}$. Für die Hyperbel empfiehlt sich besonders die Form $r = \frac{p\cos\psi}{2\cos\frac{1}{2}(v+\psi)\cos\frac{1}{2}(v-\psi)}\cdot$
